\chapter*{Abstract}
% \addcontentsline{toc}{chapter}{Abstract}
\label{chap:abstract}
% NOTE (24 Nov 2026): First draft (about 290 words) written from the Results and Conclusion chapters at the owner's request, drafted with AI assistance; reword in your own words before submission. Structure: why and aim; what was built and evaluated with the main results; what was concluded.

University timetables depend on many rules about when instructors, rooms and student groups are available. Administrators usually state these rules informally, and entering them into a timetabling system through forms is slow. This thesis investigates whether a language model can turn such requirements, written in English, formal Arabic, Egyptian Arabic or a mix of them, into formal rules for a constraint solver without allowing a wrong interpretation to reach a published timetable.

The work produced JadwAI, a web application for teaching and examination timetables. Requirements are interpreted into a closed set of rules, shown to the administrator for approval, and only approved rules are passed to a CP-SAT solver. Every generated timetable is checked by an independent validator before it can be published. The system also includes JAI, a chat assistant whose language model runs locally. On a synthetic development corpus of 144 statements, a rule-based parser produced the exact expected output for 32.6\% of statements. A local seven-billion-parameter model reached 9.0\% with the application's prompt and 16.7\% with a prompt aligned to the evaluation contract, and its intervals overlapped the parser's. The model learned to ask for missing information, but also asked when it was not needed and missed every conflict. The assistant model was fine-tuned in four rounds; guide-grounded training answers improved it in a blind rating, while one round made it worse. The solver built a validated department timetable of 448 sessions and a university-scale timetable of 2{,}703 sessions in six parts.

The results suggest that, at this scale, a local language model is useful as a proposer whose output is reviewed and validated, not as a trusted interpreter. The labels are unreviewed and the corpus is synthetic, so the findings are development evidence only; the comparison with manual entry remains future work.
